% GENERATED FILE. Edit canonical lessons, then run npm run generate:books.
% canonical-source-hash: fnv1a64:e6084aaeb04b457c
% canonical-lessons: PA-C304-khurpa, PA-C304-godna, PA-C304-pakkna, PA-C304-sitta, PA-C304-dana

\chapter{Weeding and Ripening}
\label{ch:pa-a2-weeding-and-ripening}
\hlchaptermodality{304}

\begin{chapteropening}
\textbf{By the end of this chapter:} \emph{I can talk about the crop as it grows.}

Complete the last lesson of chapter 304: I can talk about the crop as it grows.
\end{chapteropening}

\section[\texorpdfstring{\pa{ਖੁਰਪਾ} (khurpā)}{ਖੁਰਪਾ (khurpā)}]{\pa{ਖੁਰਪਾ} (khurpā) — a weeding hoe}

\label{lesson:PA-C304-khurpa}



\begin{quote}
Before the new one: say the Punjabi for a farmer, then the Punjabi for a Persian wheel.
\end{quote}

\subsection*{You'll want to know: \pa{ਖੁਰਪਾ}}
\textbf{\pa{ਖੁਰਪਾ}} — \emph{khurpā} — \textquotedblleft{}a weeding hoe\textquotedblright{}.

\begin{quote}
\pa{ਮੈਂ} \pa{ਖੁਰਪੇ} \pa{ਨਾਲ} \pa{ਘਾਹ} \pa{ਕੱਟਦਾ} \pa{ਹਾਂ।}
\end{quote}

\begin{quote}
\emph{maiṁ khurpe nāl ghāh kaṭṭdā hāṁ.}
\end{quote}

\begin{quote}
I cut the grass with a weeding hoe.
\end{quote}

A masculine noun.

\begin{cousinweb}
\textbf{\pa{ਖੁਰਪਾ}} (\emph{khurpā}) is inherited from Sanskrit \emph{kṣurapra}, \textquotedblleft{}a sort of hoe\textquotedblright{}; CDIAL lists Punjabi \emph{khurpā}, \textquotedblleft{}grubber for grass\textquotedblright{}, under it (CDIAL 3730).
\end{cousinweb}

\begin{grammarlens}[title={What you've built}]
One more word for this chapter.
\end{grammarlens}

\subsection*{Your turn}
\begin{itemize}
\raggedright
  \item \emph{Say it:} \emph{khurpā}
  \item \emph{Say it:} \emph{khurpā}, once more
  \item \emph{Say it:} \emph{maiṁ khurpe nāl ghāh kaṭṭdā hāṁ.} — I cut the grass with a weeding hoe.
  \item \emph{Recall:} say the Punjabi for a farmer, then the Punjabi for a Persian wheel, then say \emph{khurpā} again
\end{itemize}

\begin{tcolorbox}[breakable,colback=teal!4,colframe=teal!35!black,title={Before you move on}]
What does \pa{ਖੁਰਪਾ} mean? (\textquotedblleft{}A weeding hoe\textquotedblright{}.) Say it once more.
\end{tcolorbox}
\section[\texorpdfstring{\pa{ਗੋਡਣਾ} (goḍṇā)}{ਗੋਡਣਾ (goḍṇā)}]{\pa{ਗੋਡਣਾ} (goḍṇā) — to hoe; to weed with a hoe}

\label{lesson:PA-C304-godna}



\begin{quote}
Before the new one: say the Punjabi for a Persian wheel, then the Punjabi for a weeding hoe.
\end{quote}

\subsection*{You'll want to know: \pa{ਗੋਡਣਾ}}
\textbf{\pa{ਗੋਡਣਾ}} — \emph{goḍṇā} — \textquotedblleft{}to hoe; to weed with a hoe\textquotedblright{}.

\begin{quote}
\pa{ਅਸੀਂ} \pa{ਅਗਲੇ} \pa{ਹਫ਼ਤੇ} \pa{ਮੱਕੀ} \pa{ਗੋਡਾਂਗੇ।}
\end{quote}

\begin{quote}
\emph{asīṁ agle hafte makkī goḍāṅge.}
\end{quote}

\begin{quote}
We will hoe the maize next week.
\end{quote}

Past (I …): \pa{ਮੈਂ} \pa{ਕਪਾਹ} \pa{ਗੋਡੀ}, \emph{maiṁ kapāh goḍī}.

\begin{cousinweb}
\textbf{\pa{ਗੋਡਣਾ}} (\emph{goḍṇā}) is inherited from a reconstructed \emph{*khōdd}, \textquotedblleft{}dig\textquotedblright{}, in its form \emph{*gōḍḍ}; CDIAL lists Punjabi \emph{goḍṇā} and \emph{goḍḍī}, \textquotedblleft{}hoeings\textquotedblright{}, under it, and Lahnda \emph{goḍaṇ}, \textquotedblleft{}to hoe\textquotedblright{} (CDIAL 3934).
\end{cousinweb}

\begin{grammarlens}[title={What you've built}]
One more word for this chapter.
\end{grammarlens}

\subsection*{Your turn}
\begin{itemize}
\raggedright
  \item \emph{Say it:} \emph{goḍṇā}
  \item \emph{Say it:} \emph{goḍṇā}, once more
  \item \emph{Say it:} \emph{asīṁ agle hafte makkī goḍāṅge.} — We will hoe the maize next week.
  \item \emph{Recall:} say the Punjabi for a Persian wheel, then the Punjabi for a weeding hoe, then say \emph{goḍṇā} again
  \item \emph{Write it:} copy \textbf{\pa{ਗੋਡਣਾ}} once, checking the whole word
\end{itemize}

\begin{tcolorbox}[breakable,colback=teal!4,colframe=teal!35!black,title={Before you move on}]
What does \pa{ਗੋਡਣਾ} mean? (\textquotedblleft{}To hoe; to weed with a hoe\textquotedblright{}.) Say it once more.
\end{tcolorbox}
\section[\texorpdfstring{\pa{ਪੱਕਣਾ} (pakkṇā)}{ਪੱਕਣਾ (pakkṇā)}]{\pa{ਪੱਕਣਾ} (pakkṇā) — to ripen}

\label{lesson:PA-C304-pakkna}



\begin{quote}
Before the new one: say the Punjabi for a weeding hoe, then the Punjabi for to hoe.
\end{quote}

\subsection*{You'll want to know: \pa{ਪੱਕਣਾ}}
\textbf{\pa{ਪੱਕਣਾ}} — \emph{pakkṇā} — \textquotedblleft{}to ripen\textquotedblright{}.

\begin{quote}
\pa{ਵਿਸਾਖ} \pa{ਵਿੱਚ} \pa{ਕਣਕ} \pa{ਪੱਕੇਗੀ।}
\end{quote}

\begin{quote}
\emph{visākh vich kaṇak pakkegī.}
\end{quote}

\begin{quote}
The wheat will ripen in Visakh.
\end{quote}

Past (it …): \pa{ਉਹ} \pa{ਕੇਲਾ} \pa{ਪੱਕਿਆ}, \emph{oh kelā pakkiā}.

\begin{cousinweb}
\textbf{\pa{ਪੱਕਣਾ}} (\emph{pakkṇā}) is inherited from Sanskrit \emph{pakva}, \textquotedblleft{}cooked, ripe\textquotedblright{}; CDIAL lists Punjabi \emph{pakkṇā}, \textquotedblleft{}to ripen\textquotedblright{}, under it (CDIAL 7621).
\end{cousinweb}

\begin{grammarlens}[title={What you've built}]
One more word for this chapter.
\end{grammarlens}

\subsection*{Your turn}
\begin{itemize}
\raggedright
  \item \emph{Say it:} \emph{pakkṇā}
  \item \emph{Say it:} \emph{pakkṇā}, once more
  \item \emph{Say it:} \emph{visākh vich kaṇak pakkegī.} — The wheat will ripen in Visakh.
  \item \emph{Recall:} say the Punjabi for a weeding hoe, then the Punjabi for to hoe, then say \emph{pakkṇā} again
  \item \emph{Write it:} copy \textbf{\pa{ਪੱਕਣਾ}} once, checking the whole word
\end{itemize}

\begin{tcolorbox}[breakable,colback=teal!4,colframe=teal!35!black,title={Before you move on}]
What does \pa{ਪੱਕਣਾ} mean? (\textquotedblleft{}To ripen\textquotedblright{}.) Say it once more.
\end{tcolorbox}
\section[\texorpdfstring{\pa{ਸਿੱਟਾ} (siṭṭā)}{ਸਿੱਟਾ (siṭṭā)}]{\pa{ਸਿੱਟਾ} (siṭṭā) — an ear of grain}

\label{lesson:PA-C304-sitta}



\begin{quote}
Before the new one: say the Punjabi for to hoe, then the Punjabi for to ripen.
\end{quote}

\subsection*{You'll want to know: \pa{ਸਿੱਟਾ}}
\textbf{\pa{ਸਿੱਟਾ}} — \emph{siṭṭā} — \textquotedblleft{}an ear of grain\textquotedblright{}.

\begin{quote}
\pa{ਕਣਕ} \pa{ਦਾ} \pa{ਸਿੱਟਾ} \pa{ਭਾਰਾ} \pa{ਹੈ।}
\end{quote}

\begin{quote}
\emph{kaṇak dā siṭṭā bhārā hai.}
\end{quote}

\begin{quote}
The ear of wheat is heavy.
\end{quote}

A masculine noun.

\begin{cousinweb}
\textbf{\pa{ਸਿੱਟਾ}} (\emph{siṭṭā}) is inherited from a reconstructed \emph{*siṭṭa}, \textquotedblleft{}ear of corn\textquotedblright{}; CDIAL lists Punjabi \emph{siṭṭā} under it (CDIAL 13396).
\end{cousinweb}

\begin{grammarlens}[title={What you've built}]
One more word for this chapter.
\end{grammarlens}

\subsection*{Your turn}
\begin{itemize}
\raggedright
  \item \emph{Say it:} \emph{siṭṭā}
  \item \emph{Say it:} \emph{siṭṭā}, once more
  \item \emph{Say it:} \emph{kaṇak dā siṭṭā bhārā hai.} — The ear of wheat is heavy.
  \item \emph{Recall:} say the Punjabi for to hoe, then the Punjabi for to ripen, then say \emph{siṭṭā} again
  \item \emph{Write it:} copy \textbf{\pa{ਸਿੱਟਾ}} once, checking the whole word
\end{itemize}

\begin{tcolorbox}[breakable,colback=teal!4,colframe=teal!35!black,title={Before you move on}]
What does \pa{ਸਿੱਟਾ} mean? (\textquotedblleft{}An ear of grain\textquotedblright{}.) Say it once more.
\end{tcolorbox}
\section[\texorpdfstring{\pa{ਦਾਣਾ} (dāṇā)}{ਦਾਣਾ (dāṇā)}]{\pa{ਦਾਣਾ} (dāṇā) — a grain; a seed of grain}

\label{lesson:PA-C304-dana}



\begin{quote}
Before the new one: say the Punjabi for to ripen, then the Punjabi for an ear of grain.
\end{quote}

\subsection*{You'll want to know: \pa{ਦਾਣਾ}}
\textbf{\pa{ਦਾਣਾ}} — \emph{dāṇā} — \textquotedblleft{}a grain; a seed of grain\textquotedblright{}.

\begin{quote}
\pa{ਹਰ} \pa{ਸਿੱਟੇ} \pa{ਵਿੱਚ} \pa{ਬਹੁਤ} \pa{ਦਾਣੇ} \pa{ਹਨ।}
\end{quote}

\begin{quote}
\emph{har siṭṭe vich bahut dāṇe han.}
\end{quote}

\begin{quote}
There are many grains in every ear.
\end{quote}

A masculine noun.

\begin{cousinweb}
\textbf{\pa{ਦਾਣਾ}} (\emph{dāṇā}) is a loan, usually traced to Persian \emph{dāna}, \textquotedblleft{}grain, seed\textquotedblright{}.
\end{cousinweb}

\begin{grammarlens}[title={What you've built}]
One more word for this chapter.
\end{grammarlens}

\subsection*{Your turn}
\begin{itemize}
\raggedright
  \item \emph{Say it:} \emph{dāṇā}
  \item \emph{Say it:} \emph{dāṇā}, once more
  \item \emph{Say it:} \emph{har siṭṭe vich bahut dāṇe han.} — There are many grains in every ear.
  \item \emph{Recall:} say the Punjabi for to ripen, then the Punjabi for an ear of grain, then say \emph{dāṇā} again
  \item \emph{Write it:} copy \textbf{\pa{ਦਾਣਾ}} once, checking the whole word
\end{itemize}

\begin{tcolorbox}[breakable,colback=teal!4,colframe=teal!35!black,title={Before you move on}]
What does \pa{ਦਾਣਾ} mean? (\textquotedblleft{}A grain; a seed of grain\textquotedblright{}.) Say it once more.
\end{tcolorbox}
